\chapter{Smooth Muscle Antibody (SMA)}

\section*{What Is This Test?}
The smooth muscle antibody (SMA) test, also called anti-smooth muscle antibody (ASMA) or anti-actin antibody, detects autoantibodies against components of smooth muscle. The main target is F-actin (filamentous actin). SMA is the serologic hallmark of autoimmune hepatitis type 1. This is a chronic inflammatory liver disease marked by hepatocyte destruction, high serum transaminases, and hypergammaglobulinemia. The main target antigen is F-actin, a cytoskeletal protein in smooth muscle cells. SMA may also react with other cytoskeletal components, including G-actin, vimentin, desmin, and tubulin. Testing is done by indirect immunofluorescence (IIF) on composite tissue substrates (typically rodent kidney, stomach, and liver sections). This allows simultaneous detection of SMA and other liver-related autoantibodies (antimitochondrial antibody, liver-kidney microsomal antibody). High-titer SMA (>1:80) with F-actin specificity in a patient with hepatitis strongly supports a diagnosis of autoimmune hepatitis type 1.

\section*{Alternative Names}
\begin{itemize}
  \item SMA
  \item Smooth Muscle Antibody
  \item ASMA (Anti-Smooth Muscle Antibody)
  \item Anti-Actin Antibody
  \item Anti-F-Actin Antibody
  \item Anti-Smooth Muscle Antibody
\end{itemize}
\section*{Principle}
SMA is detected by indirect immunofluorescence (IIF) using composite tissue substrates made of frozen sections of rodent kidney, stomach, and liver. Patient serum is diluted (typically starting at 1:20) and incubated with the tissue substrate. After washing, fluorescein-conjugated anti-human IgG (or polyvalent anti-IgG/IgM/IgA) is applied. A trained technologist examines the stained tissue sections under a fluorescence microscope. They report both the titer (the highest dilution at which specific staining is visible) and the pattern of staining. Characteristic SMA staining shows bright fluorescence in the muscularis mucosa of the stomach, the muscularis layer of blood vessel walls in the kidney (particularly the media of small arteries and arterioles), and the perilobular and perisinusoidal areas of the liver. The VGM pattern (vessels V, glomerular mesangium G, and muscularis mucosa M) is characteristic of anti-F-actin SMA. ELISA using purified F-actin as antigen gives higher specificity for autoimmune hepatitis type 1 than IIF alone.

\section*{History}
SMA was first described in 1965 by Johnson and colleagues in patients with chronic active hepatitis. The link between high-titer SMA and autoimmune hepatitis was firmly established in the 1970s. In the 1980s, Kurki et al. and Toh et al. identified F-actin as the main target antigen of SMA in autoimmune hepatitis. The 1999 International Autoimmune Hepatitis Group (IAIHG) diagnostic criteria added SMA as one of the immunologic markers for autoimmune hepatitis. The simplified IAIHG scoring system, published in 2008, gives points for the presence of autoantibodies (ANA, SMA, anti-LKM-1, or anti-LC-1) in the diagnosis of autoimmune hepatitis. Autoimmune hepatitis was split into type 1 (ANA and/or SMA positive) and type 2 (anti-LKM-1 and/or anti-LC-1 positive) based on distinct serologic profiles.

\section*{Why This Test Is Ordered}
\begin{itemize}
  \item Diagnosing autoimmune hepatitis type 1 in patients with acute or chronic hepatitis of unknown cause, high serum transaminases (AST and ALT), and hypergammaglobulinemia (especially high IgG)
  \item Distinguishing autoimmune hepatitis from viral hepatitis (hepatitis A, B, C, EBV, CMV), drug-induced liver injury, Wilson disease, alpha-1 antitrypsin deficiency, and non-alcoholic steatohepatitis (NASH)
  \item Checking unexplained high liver enzymes in a patient with signs of autoimmune disease (arthralgias, rash, hemolytic anemia) or other autoimmune conditions (autoimmune thyroiditis, celiac disease, inflammatory bowel disease)
  \item Classifying autoimmune hepatitis into type 1 (SMA and/or ANA positive, affecting mostly adults and adolescents) or type 2 (anti-LKM-1 positive, affecting mostly children)
  \item Monitoring disease activity and response to immunosuppressive therapy in known autoimmune hepatitis, though SMA titer does not track closely with histologic activity
  \item Pre-treatment evaluation before corticosteroid or azathioprine therapy for autoimmune hepatitis
\end{itemize}

\section*{Primary vs Secondary Test}
This is a primary diagnostic test for autoimmune hepatitis type 1. SMA, along with ANA, is a first-line autoantibody screening test in the evaluation of unexplained hepatitis. A positive high-titer SMA strongly supports the diagnosis. However, liver biopsy must confirm it.

\section*{How This Test Is Performed}
For most patients, a health care professional takes a blood sample from a vein in your arm using a small needle. After the needle goes in, a small amount of blood is collected into a test tube or vial. You may feel a little sting when the needle goes in or out. This usually takes less than five minutes.

\section*{What Preparation Is Needed}
No special preparation is required. Fasting is not necessary. Patients should tell their healthcare provider about all medications, because certain drugs (methyldopa, nitrofurantoin, minocycline, anti-TNF agents, statins) can cause drug-induced autoimmune-like hepatitis with positive SMA. Testing is ideally done before starting immunosuppressive therapy (corticosteroids, azathioprine), because treatment may lower autoantibody titers.

\section*{Contraindications and Precautions}
\begin{itemize}
  \item There are no absolute contraindications. The specimen should be collected in a serum separator tube (SST, gold-top) or red-top tube. Hemolyzed, lipemic, or icteric specimens may interfere with immunofluorescence interpretation. Low-titer SMA (<1:40) is common in the general population (up to 5\% of healthy adults). It is also found in non-autoimmune liver disease (viral hepatitis B and C, alcoholic liver disease, non-alcoholic fatty liver disease). SMA is also detected in a large share of patients with other autoimmune conditions: celiac disease (30--50\%), primary biliary cholangitis, and chronic HCV infection. F-actin-specific SMA (VGM pattern by IIF or positive anti-F-actin ELISA) has better specificity (>90\%) for autoimmune hepatitis type 1 than non-actin SMA. ANA and SMA testing should be done together when autoimmune hepatitis is suspected, because about 20--30\% of patients have only one of these antibodies.
\end{itemize}
\section*{Risks}
The risk is minimal. Slight pain or bruising at the venipuncture site. Rare: hematoma, infection, vasovagal syncope.

\section*{What Happens After the Test}
Results are usually available within 2 to 5 days. A positive SMA result should be read alongside ANA, anti-LKM-1, serum IgG level, liver enzymes (AST, ALT, alkaline phosphatase), and liver biopsy findings. High-titer SMA (>1:80) with a compatible clinical and biochemical profile strongly supports autoimmune hepatitis type 1. Liver biopsy is essential. It confirms the diagnosis, assesses histologic grade (inflammation) and stage (fibrosis), and excludes other liver diseases. Serial SMA testing is of limited value for monitoring disease activity. Serum transaminases (ALT, AST) and IgG levels are the preferred laboratory markers for monitoring treatment response.

\section*{Understanding Results}

\subsection*{Below Normal (Negative)}
\begin{itemize}
  \item A negative SMA does not exclude autoimmune hepatitis type 1 --- about 20--30\% of patients with biopsy-proven autoimmune hepatitis type 1 are SMA-negative but ANA-positive
  \item About 5--10\% of patients with autoimmune hepatitis type 1 are negative for both SMA and ANA. Diagnosis then relies on high serum IgG, characteristic liver histology, and excluding other causes
  \item Anti-LKM-1 (anti-liver-kidney microsomal type 1) or anti-LC-1 (anti-liver cytosol type 1) positive autoimmune hepatitis (type 2) is SMA-negative and ANA-negative
  \item Acute severe autoimmune hepatitis presenting as acute liver failure may be seronegative at first presentation
  \item Viral hepatitis (hepatitis A, B, C, EBV, CMV) is usually SMA-negative, though transient low-titer SMA may be seen
  \item Non-alcoholic fatty liver disease (NAFLD) and non-alcoholic steatohepatitis (NASH) are usually SMA-negative
  \item Drug-induced liver injury (acetaminophen, antibiotics, NSAIDs, herbal supplements) is usually SMA-negative unless a specific drug triggers autoimmune-like hepatitis (nitrofurantoin, minocycline, anti-TNF agents)
\end{itemize}

\subsection*{Above Normal (Positive)}
\begin{itemize}
  \item Autoimmune hepatitis type 1: SMA is positive in 70--80\% of patients with autoimmune hepatitis type 1. High-titer SMA ($\geq$1:80) with F-actin specificity (VGM pattern on tissue substrate) is highly specific for autoimmune hepatitis type 1, especially in patients with high serum transaminases and hypergammaglobulinemia (IgG >1.5 times upper limit of normal)
  \item ANA/SMA double positivity: Found in about 50--60\% of patients with autoimmune hepatitis type 1. Double positivity strengthens the serologic evidence for autoimmune hepatitis
  \item Celiac disease: SMA (non-actin specificity) is positive in 30--50\% of patients with celiac disease consuming gluten. The SMA usually disappears with gluten-free diet adherence and is not linked to liver disease
  \item Viral hepatitis: Low-titer SMA (<1:40) is found in 10--20\% of patients with chronic viral hepatitis B or C from non-specific polyclonal B-cell activation. Anti-F-actin ELISA is usually negative
  \item Primary biliary cholangitis (PBC): SMA is positive in a minority of patients with PBC, but antimitochondrial antibody (AMA) positivity is the hallmark serologic finding. AMA-positive PBC with concurrent SMA positivity represents an autoimmune hepatitis-PBC overlap syndrome
  \item Drug-induced autoimmune-like hepatitis: Nitrofurantoin, minocycline, anti-TNF biologics (infliximab, adalimumab), and statins can cause a syndrome of liver injury with positive ANA and SMA that mimics autoimmune hepatitis. Stopping the drug usually leads to resolution, unlike idiopathic autoimmune hepatitis, which usually requires long-term immunosuppression
  \item Acute severe autoimmune hepatitis: Markedly high SMA titers ($\geq$1:320) may be present at presentation. However, a negative SMA does not exclude the diagnosis, and a liver biopsy should not be delayed in patients with suspected acute severe autoimmune hepatitis
\end{itemize}

\section*{Normal Results}
\begin{itemize}
  \item \textbf{SMA by indirect immunofluorescence:} Negative (no specific smooth muscle staining at a 1:20 dilution)
  \item \textbf{Borderline:} Titer 1:20 to 1:40 (equivocal; clinical significance is low in the absence of compatible clinical and biochemical findings)
  \item \textbf{Low-titer positive:} Titer 1:40 (may be seen in 5--10\% of healthy adults and in various non-autoimmune liver diseases)
  \item \textbf{Significant titer:} Titer $\geq$1:80 (clinically significant, particularly with F-actin specificity and compatible clinical features)
  \item \textbf{Anti-F-actin ELISA:} <20 U/mL (negative)
  \item F-actin-specific SMA (positive anti-F-actin ELISA or VGM pattern on tissue substrate) has >90\% specificity for autoimmune hepatitis type 1
  \item SMA results should always be read alongside ANA, anti-LKM-1, serum IgG levels, liver enzymes (AST, ALT, alkaline phosphatase), and liver biopsy findings
\end{itemize}

\section*{Follow-up Tests}
\begin{itemize}
  \item \textbf{Antinuclear antibody (ANA):} Should be ordered at the same time as SMA for evaluation of autoimmune hepatitis. ANA is positive in 50--70\% of patients with autoimmune hepatitis type 1
  \item \textbf{Anti-LKM-1 (anti-liver-kidney microsomal type 1) and anti-LC-1 (anti-liver cytosol type 1):} To diagnose autoimmune hepatitis type 2, which is SMA-negative and ANA-negative. Primarily seen in children and young adults
  \item \textbf{Anti-F-actin ELISA:} Confirmatory test for F-actin specificity. Higher specificity for autoimmune hepatitis type 1 than IIF SMA alone
  \item \textbf{Serum IgG:} High IgG (>1.5 times upper limit of normal) is characteristic of autoimmune hepatitis and is included in the simplified IAIHG diagnostic criteria
  \item \textbf{Serum AST, ALT, alkaline phosphatase, total bilirubin:} Liver function tests to assess the degree of hepatocellular injury (AST, ALT) and cholestasis (alkaline phosphatase, bilirubin)
  \item \textbf{Liver biopsy:} Essential for diagnosis. Characteristic histologic findings include interface hepatitis (piecemeal necrosis) with lymphoplasmacytic infiltration, hepatocyte rosetting, and emperipolesis (intact lymphocytes within hepatocyte cytoplasm)
  \item \textbf{Viral hepatitis serology (hepatitis A IgM, HBsAg, HBc IgM, anti-HCV, EBV, CMV):} To exclude viral hepatitis
  \item \textbf{Antimitochondrial antibody (AMA):} To exclude primary biliary cholangitis, which is characterized by AMA positivity (>95\%) and different histologic findings (florid bile duct lesions)
  \item \textbf{Serum ceruloplasmin and 24-hour urinary copper:} To exclude Wilson disease in young patients (<40 years) with unexplained hepatitis
  \item \textbf{Serum alpha-1 antitrypsin level and phenotype:} To exclude alpha-1 antitrypsin deficiency
  \item \textbf{Celiac disease serology (tissue transglutaminase IgA, total IgA):} If celiac disease is suspected (diarrhea, weight loss, iron deficiency anemia) in a patient with positive low-titer SMA
\end{itemize}

\section*{References}
\begin{enumerate}
  \item Hennes EM, Zeniya M, Czaja AJ, et al. Simplified criteria for the diagnosis of autoimmune hepatitis. Hepatology. 2008;48(1):169--176.
  \item Manns MP, Czaja AJ, Gorham JD, et al. Diagnosis and management of autoimmune hepatitis. Hepatology. 2010;51(6):2193--2213.
  \item Czaja AJ. Autoantibodies in autoimmune liver disease. Adv Clin Chem. 2005;40:127--164.
  \item Krawitt EL. Autoimmune hepatitis. N Engl J Med. 2006;354(1):54--66.
  \item Verdonk RC, Lozano MF, van den Berg AP, et al. The significance of liver biopsy in the diagnostic approach of autoimmune hepatitis. Neth J Med. 2011;69(9):387--393.
\end{enumerate}

\clearpage
\section*{Regulatory Status and Authorization Date}
This chapter describes a test category, not a specific commercial product. FDA, CDSCO, EU IVDR, and MHRA approval or registration dates therefore cannot be assigned without the assay manufacturer, product name, instrument, and intended use. For laboratory-developed tests, an individual agency approval date may not exist. See the Preface for the interpretation of regulatory status.

\clearpage
